\documentclass[10pt,a4paper]{amsart}
\usepackage[margin=19mm]{geometry}
\usepackage{amsmath,amssymb,mathtools}
\usepackage[scr=rsfs,cal=euler]{mathalfa}
\usepackage{xfrac}
\usepackage[unicode,hidelinks,bookmarks=false]{hyperref}
\newcommand{\ie}{i.e.\ }
\newcommand{\cf}{cf.\ }
\newcommand{\resp}{respectively\ }
\newcommand{\Subsection}{\subsection}
\providecommand{\nobreakdash}{\nobreak\hbox{-}}
\setlength{\emergencystretch}{2em}
\multlinegap0pt
\usepackage{bm}
\usepackage{amssymb}
\usepackage{tikz}
\usepackage[shortlabels]{enumitem}
\usepackage{float}
\newtheorem{prop}{Proposition}
\newtheorem{thm}{Theorem}
\def\cA{\mathcal{A}}
\def\cH{\mathcal{H}}
\def\cI{\mathcal{I}}
\def\dd{\mathrm{d}}
\def\frc{\mathfrak{F}}
\def\frf{\mathfrak{f}}
\def\fru{\mathfrak{u}}
\def\frv{\mathfrak{v}}
\title{Complex spinorial forms, Brinkmann four-manifolds, and self-dual bundle gerbes}
\author{Alejandro Gil-Garc\'ia, C. S. Shahbazi}
\date{Source: arXiv:2605.22798v1 (CC BY 4.0)}
\begin{document}
\maketitle

\noindent\emph{Source and licence.} Complex spinorial forms, Brinkmann four-manifolds, and self-dual bundle gerbes. Version arXiv:2605.22798v1 (CC BY 4.0), distributed by arXiv under the Creative Commons Attribution 4.0 International licence, http://creativecommons.org/licenses/by/4.0/, as recorded in the arXiv licence metadata for this exact version and shown on the abstract page https://arxiv.org/abs/2605.22798v1. Full licence notice: \texttt{licenses/arxiv-2605.22798-CC-BY-4.0.txt}.\par\smallskip

\noindent\textit{Editorial scope.} Counted unit: the article's thm thm:NoGoFreedman (source0.tex lines 2278--2291). The article's complete original proof of it is delivered with it (source0.tex lines 2293--2363, one \texttt{proof} environment of 71 lines). No mathematics is written, repaired or re-derived: the statement and the proof are the article's own lines, in the article's own notation, and no retained line is substituted or re-flowed.  Retained uncounted as the source-local context the proof needs: lines 1805--1808; lines 2026--2029; lines 2064--2071; lines 2229--2239.  Nothing was omitted from any retained span, so the package carries no omission record.  No retained line contains a citation, so the wrapper carries no bibliography and no bibliography entry.  No retained line contains a figure environment, an image-inclusion command or drawing code, so no image is delivered and the package declares no asset dependency.  Editorial formatting only, in addition: an AMS \texttt{amsart} preamble that restates the article's own declarations and macros used by the retained lines, namely the environments \texttt{prop}, \texttt{thm} on separate counters, together with the standard packages those lines need.  The retained environments print in this document as Proposition 1, Theorem 1 in order of appearance, whereas the article prints the same items with the numbering of its own full document.  The complete native source of the article as distributed in the arXiv e-print for this exact version is bundled unchanged as \texttt{sources/201161/source0.tex}.
\begin{equation}
\label{eq:Freedmaneqs}
\mathrm{Ric}^g - \frac{1}{2} s_g g  = \Lambda g + \mathfrak{e}^2\, F_A \circ_g F_A - \frac{\mathfrak{e}^2}{2}  \vert F_A \vert^2_g g\, , \qquad \dd \ast_g F_A =0 
\end{equation}

\begin{equation}
\label{eq:standardBrinkmann}
g = \cH_{\fru} \dd \fru \otimes \dd \fru +  \dd \fru \odot (\dd \frv + \cA_{\fru}) + h_{\fru}\, , \qquad  u = \dd \fru\, , \qquad \partial_{\frv}A = 0,
\end{equation}

\begin{prop}
\label{prop:brinkmannsusy}
A standard Brinkmann configuration $(h_{\fru}, a_{\fru}, \cA_{\fru}, \psi_{\fru}, \phi_{\fru}, \cH_{\fru})$ on $(X,P)$ is quasi-supersymmetric if and only if it satisfies the following differential system:
\begin{equation*}
\phi_{\fru} = \lambda_R \fru + \lambda_R^o \, , \qquad \ast_{h_{\fru}} F_{a_{\fru}} = \mu \lambda_I,
\end{equation*}
where $\lambda_R^o\in \mathbb{R}$ is a real constant.
\end{prop}

\begin{prop}
\label{prop:equivalentBrinkmann}
A quasi-supersymmetric stationary configuration  $(h, a, \cA, \psi, \phi, \cH)$ on $(X,P)$ satisfies Freedman's gauged supergravity equations \eqref{eq:Freedmaneqs} on $\cI_u\times \cI_v\times X$ if and only if:
\begin{eqnarray*}
& \vert F_{\cA}\vert^2_{h} + \nabla^{h\ast}\dd \cH     =   2 \mathfrak{e}^2 \vert \dd \psi\vert^2_h \, , \qquad \nabla^{h\ast} F_{\cA}   =     2 \mathfrak{e}^2 \mu \lambda_I \ast_h \dd\psi \, , \qquad \dd\ast_h \dd \psi   = \mu\lambda_I F_{\cA}\, ,\\
&   s_h = 2 \mathfrak{e}^2 \lambda_I^2\, , \qquad    2 \Lambda  +  \mathfrak{e}^2 \lambda_I^2 = 0\, ,
\end{eqnarray*}

\noindent
where $g$ is given as in Equation \eqref{eq:standardBrinkmann} in terms of the given tuple $(h, a, \cA, \psi, \phi, \cH)$ and $F_{\cA} = \dd\cA$.
\end{prop}

\begin{thm}
\label{thm:NoGoFreedman}
A standard stationary configuration  $(h, a, \cA, \psi, \phi, \cH)$ on the pair $(X,P)$ is a quasi-supersymmetric solution of Freedman's gauged supergravity equations \eqref{eq:Freedmaneqs} on $\mathbb{R}^2 \times X$ if and only if all the following conditions hold:
\begin{itemize}
    \item $X = S^2$, that is, $X$ is diffeomorphic to the two-dimensional sphere.
    \item $h$ is the round metric of scalar curvature $s_h = 2 \mathfrak{e}^2 \lambda_I^2$, that is, of radius $R=\vert\mathfrak{e} \lambda_I\vert^{-1}$.
    \item The Chern number $c(P)$ satisfies $\lambda_I \mathfrak{e}^2 c(P) = 2\mu$, and $a$ is a Yang-Mills connection on $P$.  
    \item $\psi \in C^{\infty}(S^2)$ is an eigenfunction of the Laplacian on $(S^2,h)$ with the lowest non-zero eigenvalue. 
    \item $\cA \in \Omega^1(S^2)$ is given by $\cA = \mu \lambda^{-1}_I (\ast_h \dd \psi + \dd \mathfrak{f})$ for a function $\mathfrak{f}\in C^{\infty}(S^2)$.
    \item $\cH\in C^{\infty}(S^2)$ is given by $\lambda_I^2\cH = -(R^{-2}\psi^2 + \frc)$ for a real constant $\frc$.
    \item The cosmological constant $\Lambda$ is determined by $2\Lambda = -\mathfrak{e}^2 \lambda_I^2$ and $\phi$ is constant.
\end{itemize}
In particular, the aforementioned Poisson equation for $\cH$ is guaranteed to have a solution.
\end{thm}

\begin{proof}
Let $(h, a, \cA, \psi, \phi, \cH)$ be a stationary configuration of Freedman's gauged supergravity on $\mathbb{R}^2 \times X$. By Proposition \ref{prop:brinkmannsusy}, the quasi-supersymmetry conditions reduce to $\lambda_R = 0$, hence $\phi$ is constant, and $a$ being a Yang-Mills connection on $P$ with normalization:
\begin{equation*}
F_a = \mu\lambda_I \nu_h.
\end{equation*}

\noindent
Equivalently, the Chern number $c(P)$ of $P$ satisfies:
\begin{equation*}
c(P) = \frac{\mu \chi(X)}{\lambda_I \mathfrak{e}^2},
\end{equation*}

\noindent
where $\chi(X)$ is the Euler characteristic of $X$. These conditions, together with Proposition \ref{prop:equivalentBrinkmann}, give the necessary and sufficient conditions for the tuple $(h, a, \cA, \psi, \phi, \cH)$ being a stationary quasi-supersymmetric solution of Freedman's gauged supergravity equations \eqref{eq:Freedmaneqs}. By condition $s_h = 2\mathfrak{e}^2 \lambda_I^2$ in Proposition \ref{prop:equivalentBrinkmann} it immediately follows that $(X,h)$ is the round sphere of radius $R := \vert \mathfrak{e} \lambda_I \vert^{-1}$ since $h$ is complete by definition. Hence $\chi(X)=2$. On the other hand, the equation $\dd\ast_h \dd \psi   = \mu\lambda_I F_{\cA}$ in Proposition \ref{prop:equivalentBrinkmann} can be explicitly integrated by setting:
\begin{equation*}
\cA = \mu \lambda^{-1}_I (\ast_h \dd \psi + \dd \mathfrak{f})
\end{equation*}

\noindent
for a smooth function $\frf \in C^{\infty}(X)$. Substituting in the remaining conditions contained in Proposition \ref{prop:equivalentBrinkmann}, and setting $2 \Lambda  := -  \mathfrak{e}^2 \lambda_I^2$, we obtain:
\begin{equation*}
\lambda_I^{-2}\vert \dd\ast_h \dd\psi \vert^2_{h} + \nabla^{h\ast}\dd \cH     =   2 \mathfrak{e}^2 \vert \dd \psi\vert^2_h \, , \qquad \nabla^{h\ast} \dd \ast_h \dd\psi   =     2 \mathfrak{e}^2   \lambda^2_I \ast_h \dd\psi  
\end{equation*}

\noindent
as the remaining system that needs to be solved. We observe that the second equation above can be rewritten as follows:
\begin{equation*}
\nabla^{h\ast} \dd \ast_h \dd\psi  =  \dd^{h\ast}\dd \ast_h \dd \psi =    (\dd^{h\ast}\dd + \dd \dd^{h\ast}) \ast_h \dd \psi = \Delta_h \ast_h \dd \psi = 2 \mathfrak{e}^2   \lambda^2_I \ast_h \dd\psi = \frac{2}{R^2} \ast_h \dd\psi,
\end{equation*}

\noindent
where $\dd^{h\ast}$ is the formal adjoint of the exterior derivative, $R$ is the radius of $(S^2,h)$, and $\Delta_h$ is the positive Laplace operator on $(S^2,h)$. Hence, we conclude that $\psi$ is, modulo a constant, a linear combination of the three linearly independent eigenfunctions of the Laplacian with the lowest eigenvalue. Using standard notation for the spherical harmonics $Y^m_l$ on $(S^2,h)$, we have: 
\begin{equation*}
\psi = \sum_{m =-1}^1 Y_1^m\, c_m
\end{equation*}

\noindent
for constants $c_1 , c_0, c_{-1}$. It remains to solve:
\begin{equation*}
\nabla^{h\ast}\dd \cH     =   2 \mathfrak{e}^2 \vert \dd \psi\vert^2_h   - \lambda_I^{-2}\vert \dd\ast_h \dd\psi \vert^2_{h}.
\end{equation*}

\noindent
Integrating, we obtain:
\begin{equation*}
\int_{S^2} \vert \dd\ast_h \dd\psi \vert^2_{h} \nu_h = \int_{S^2} \langle  \dd^{h\ast}\dd\psi , \dd^{h\ast}\dd\psi \rangle_{h} = \int_{S^2} \langle   \dd\psi , \dd \dd^{h\ast}\dd\psi \rangle_{h}= \int_{S^2} \langle  \ast_h  \dd\psi , \Delta_h \ast_h\dd\psi \rangle_{h} = 2 \mathfrak{e}^2 \lambda_I^2  \vert \dd \psi\vert_h^2.
\end{equation*}

\noindent
Hence:
\begin{equation*}
\int_{S^2} \nabla^{h\ast}\dd \cH  \nu_h   =  \int_{S^2}    (2 \mathfrak{e}^2 \vert \dd \psi\vert^2_h   - \lambda_I^{-2}\vert \dd\ast_h \dd\psi \vert^2_{h}) \nu_h  =  \int_{S^2}   (2\mathfrak{e}^2 \vert \dd \psi\vert^2_h   - 2 \mathfrak{e}^2   \vert \dd \psi\vert_h^2 ) \nu_h = 0
\end{equation*}
and thus, crucially, the standard theory of elliptic operators implies the existence of a solution $\cH$ that is unique up to an additive constant. Furthermore, we have:
\begin{equation*}
\lambda_I^{-2}\vert \dd\ast_h \dd\psi \vert^2_{h} = \lambda_I^{-2} (\Delta_h \psi)^2 = \frac{4}{R^4 \lambda^2_I} \psi^2.
\end{equation*}

\noindent
Hence, the Poisson equation for $\psi$ reduces to:
\begin{equation*}
\lambda^2_I \nabla^{h\ast}\dd \cH     =   \frac{2}{R^2} \vert \dd \psi\vert^2_h   - \frac{4}{R^4}\psi^2.
\end{equation*}

\noindent
Now, we observe that the following choice:
\begin{equation*}
\lambda_I^2\cH=-(R^{-2}\psi^2+\frc),
\end{equation*}
where $\frc$ is a real constant, is a solution of the Poisson equation for $\cH$. Since solutions to such an equation are unique up to an additive constant, this is the general solution.
\end{proof}

\end{document}
